% TOP_PROBLEM: {"problemId": "problem.exact-real-grothendieck-constant", "canonicalTitle": "Exact Real Grothendieck Constant", "releaseRank": 310, "primaryDomain": "analysis_pde", "primaryDomainLabel": "Analysis & PDE", "displayStatus": "open", "releaseStatus": "open"}
% !TeX program = lualatex
% Definition and sources only; this file is not an LLM solution attempt.
\documentclass[11pt]{article}
\usepackage[margin=25mm]{geometry}
\usepackage{fontspec}
\setmainfont{Latin Modern Roman}
\usepackage{amsmath,amssymb,mathtools}
\usepackage{unicode-math}
\setmathfont{Latin Modern Math}
\usepackage{xurl,hyperref}
\hypersetup{colorlinks=true,urlcolor=blue}
\setcounter{secnumdepth}{0}
\setlength{\parindent}{0pt}
\setlength{\parskip}{5pt}
\setlength{\emergencystretch}{3em}
\begin{document}
\section*{\texorpdfstring{Exact Real Grothendieck Constant}{Exact Real Grothendieck Constant}}\label{310}
\subsection{Definitions and mathematical statement}
For a real $m\times n$ matrix $A=(a_{ij})$, put
\[
 B(A)=\max_{{\varepsilon}_i,\delta_j\in\{-1,1\}}
 \left|\sum_{i,j}a_{ij}{\varepsilon}_i\delta_j\right|.
\]
The real Grothendieck constant is the least $K$ such that, for every $m,n,A$ and all unit vectors $u_i,v_j$ in any real Hilbert space,
\[
 \left|\sum_{i,j}a_{ij}\langle u_i,v_j\rangle\right|\le K B(A).
\]
Determine this optimal universal constant $K_G^{{\mathbb{R}}}$ exactly. The dimensions and the matrix are arbitrary, and the comparison is with independent signs on rows and columns. Neither a fixed-dimensional variant nor the complex Grothendieck constant is the selected quantity. An improved upper or lower estimate alone is not an exact determination.
\par\smallskip\textit{Formulation and concept references:} \href{https://arxiv.org/abs/1103.6161}{[S1]}.

\subsection{Short English statement}
What is the exact largest advantage that real unit-vector inner products can have over signs in a bilinear form?
\subsection{Sources}
\begin{itemize}
\item[S1]1103.6161v3, The Grothendieck constant is strictly smaller than Krivine's bound.
\url{https://arxiv.org/abs/1103.6161}.
\textit{Formulation source.}
\end{itemize}
\end{document}
